\section{Boundary zeros and finite-volume rigidity}\label{sec:boundary-zeros}

In this section we show that $y = -1$ is a uniform double zero of the partition-function polynomials $Z_m(y)$ for all $m \ge 2$, a structural rigidity that goes beyond the identity $Z(t,-1) \equiv 1$.

\begin{proposition}[Uniform double-zero boundary at $y = -1$]\label{prop:double-zero}
Let $u = y + 1$.
Then:
\begin{enumerate}
\item $Z_0(-1) = 1$ and $Z_1(y) = y + 1 = u$;
\item for all $m \ge 2$, the polynomial $Z_m(y)$ is divisible by $(y+1)^2$, i.e.,
\begin{equation}\label{eq:double-zero}
(y+1)^2 \mid Z_m(y) \qquad (m \ge 2).
\end{equation}
\end{enumerate}
In particular,
\[
Z_m(-1) = 0 \quad (m \ge 1), \qquad
\partial_y Z_m(-1) = 0 \quad (m \ge 2).
\]
\end{proposition}

\begin{proof}
We establish the factorization at the level of the generating function.
Writing $u = y + 1$ and $y = u - 1$, direct algebraic manipulation yields the identity
\begin{equation}\label{eq:N-D-identity}
N(t,y) = (1 + ut)\,D(t,y) + u^2\,t^2\,(1 - t^2)(1 + ty).
\end{equation}
Dividing both sides by $D(t,y)$:
\begin{equation}\label{eq:Z-factored}
Z(t,y) = \frac{N(t,y)}{D(t,y)}
= 1 + (y+1)\,t + (y+1)^2\,t^2 \cdot \frac{(1 - t^2)(1 + ty)}{D(t,y)}.
\end{equation}
The power-series expansion of the right-hand side in $t$ shows:
\begin{itemize}
\item the constant term is $1$, confirming $Z_0(y) = 1$;
\item the coefficient of $t$ is $y + 1$, confirming $Z_1(y) = y + 1$;
\item for $m \ge 2$, each coefficient $Z_m(y)$ receives a factor $(y+1)^2$ from the prefactor of $t^2$ in the last term, establishing \eqref{eq:double-zero}.
\end{itemize}
The evaluation $Z_m(-1) = 0$ for $m \ge 1$ follows from $(y+1) \mid Z_m(y)$ (which holds for all $m \ge 1$).
The vanishing of $\partial_y Z_m(-1)$ for $m \ge 2$ follows from the stronger divisibility $(y+1)^2 \mid Z_m(y)$.
\end{proof}

\begin{remark}\label{rem:boundary-rigidity}
The identity $Z(t,-1) \equiv 1$, which appears in the original analysis as a degeneracy condition (Appendix~\ref{app:degeneracies}), is now revealed as the leading symptom of a deeper structural property: the entire finite-volume sequence $\{Z_m(y)\}_{m \ge 2}$ vanishes to second order at $y = -1$.
This uniform double-zero boundary constrains the zero-locus geometry of $Z_m(y)$ near $y = -1$ for all system sizes simultaneously, and may be interpreted as a form of spectral rigidity: the point $y = -1$ is not merely a degeneracy of the generating function, but an algebraically enforced boundary of the partition-function family.
\end{remark}

